% !TeX program = xelatex
% !TeX encoding = UTF-8
% 编译：xelatex -interaction=nonstopmode web_report01.tex（运行两遍）
% 依赖：ctex、pgfplots 等常用 TeX Live / MiKTeX 宏包；图表无需外部图片。
% 基本信息中的空白项须在提交前填写。
\documentclass[UTF8,a4paper,11pt,fontset=fandol]{ctexart}
\usepackage[left=24mm,right=24mm,top=24mm,bottom=24mm,headheight=15pt,headsep=8mm]{geometry}
\usepackage{fontspec}
% 优先使用常见系统字体；缺失时使用 TeX 发行版附带字体。
\IfFontExistsTF{Times New Roman}{\setmainfont{Times New Roman}}{\setmainfont{Latin Modern Roman}}
\IfFontExistsTF{Arial}{\setsansfont{Arial}}{\setsansfont{Latin Modern Sans}}
\setmonofont{Latin Modern Mono}[Scale=MatchLowercase]
\IfFontExistsTF{SimSun}{\setCJKmainfont{SimSun}[AutoFakeBold=2.5]}{}
\IfFontExistsTF{SimHei}{\setCJKsansfont{SimHei}}{}
\usepackage{amsmath,amssymb,mathtools}
\usepackage{booktabs,tabularx,array}
\usepackage{graphicx}
\usepackage[font=small,labelfont=bf,labelsep=quad]{caption}
\usepackage{enumitem}
\usepackage{pgfplots}
\pgfplotsset{compat=1.18}
\usepackage{fancyhdr}
\usepackage{float,needspace}

\usepackage{xurl}
\usepackage[unicode,hidelinks]{hyperref}
\usepackage{bookmark}
\usepackage{microtype}
\usepackage{lastpage}
\hypersetup{pdftitle={Web 信息处理与应用：实验一——信息检索},pdfauthor={龚夏瑜、黄蕴、陈熙},pdfsubject={Qilin-IR-50K 分词、倒排索引、索引压缩与检索实验}}
\setlength{\parindent}{2em}
\setlength{\parskip}{2pt}
\linespread{1.16}
\setlength{\emergencystretch}{2em}
\setlength{\textfloatsep}{12pt plus 3pt minus 2pt}
\setlength{\intextsep}{10pt plus 2pt minus 2pt}
\setlength{\floatsep}{10pt plus 2pt minus 2pt}
\setlength{\tabcolsep}{6pt}
\renewcommand{\arraystretch}{1.18}
\setlist{leftmargin=2em,itemsep=2pt,topsep=4pt,parsep=0pt}
\setlist[enumerate]{label=\arabic*.}
\ctexset{
  section={format=\large\sffamily\bfseries,beforeskip=16pt,afterskip=8pt},
  subsection={format=\normalsize\sffamily\bfseries,beforeskip=11pt,afterskip=5pt},
  subsubsection={format=\normalsize\sffamily\bfseries,beforeskip=9pt,afterskip=4pt}
}
\setcounter{secnumdepth}{3}
\renewcommand{\figurename}{图}
\renewcommand{\tablename}{表}
\captionsetup{skip=5pt}
\captionsetup[table]{position=top}
\urlstyle{tt}
\pagestyle{fancy}
\fancyhf{}
\fancyhead[L]{\small Web 信息处理与应用}
\fancyhead[R]{\small 实验一：信息检索}
\fancyfoot[C]{\small 第\,\thepage\,页\quad 共\,\pageref{LastPage}\,页}
\renewcommand{\headrulewidth}{0.4pt}
\renewcommand{\footrulewidth}{0pt}
\raggedbottom
\clubpenalty=10000
\widowpenalty=10000
\displaywidowpenalty=10000
\newcommand{\term}[1]{#1}
\newcommand{\method}[1]{#1}
\newcommand{\tablesource}[1]{\par\vspace{3pt}\noindent\parbox{\linewidth}{\footnotesize 数据来源：\path{#1}。}}
\newcommand{\tablenote}[1]{\par\vspace{3pt}\noindent\parbox{\linewidth}{\footnotesize 注：#1}}
\newcolumntype{Y}{>{\raggedright\arraybackslash}X}

% ====== 基本信息：提交前填写 ======
\newcommand{\leadername}{\underline{\hspace{20mm}}}
\newcommand{\leaderid}{\underline{\hspace{30mm}}}
\newcommand{\members}{\begin{tabular}[t]{@{}ll@{}}龚夏瑜 & PB24331921\\黄蕴 & PB24331918\\陈熙 & PB24331924\end{tabular}}
\newcommand{\cloudlink}{\underline{\hspace{70mm}}}
\newcommand{\sysinfo}{Windows 11；Intel i7-13650HX；16\,GB 内存；Python 3.13}
% ================================

\begin{document}
\pagenumbering{roman}
\begin{titlepage}
  \thispagestyle{empty}
  \centering
  \vspace*{8mm}
  {\large Web 信息处理与应用}\par
  \vspace{7mm}
  {\sffamily\bfseries\fontsize{23}{31}\selectfont 实验一\quad 信息检索}\par
  \vspace{4mm}
  {\normalsize Qilin-IR-50K 中文信息检索实验报告}\par
  \vspace{16mm}
  \begingroup
  \renewcommand{\arraystretch}{1.55}
  \begin{tabularx}{0.94\linewidth}{@{}p{25mm}Y@{}}
    组长 & \leadername\quad 学号：\leaderid\\[3pt]
    小组成员 & \members\\[4pt]
    运行环境 & \sysinfo\\[3pt]
    附件索引 & \cloudlink\\[3pt]
    报告日期 & 2026 年 10 月
  \end{tabularx}
  \endgroup
  \vfill
  \begin{minipage}{0.94\linewidth}
    \begin{center}\bfseries 摘\quad 要\end{center}
    \small
    本实验以 Qilin-IR-50K 的 50,000 篇中文笔记为对象，比较分词方法，构建文档级和位置级倒排索引，并测试索引压缩、布尔检索、TF-IDF 与语义检索。FMM 保留停用词后得到 115,215 个词项和 4,074,930 条词项--文档关系。Front Coding 与 Gap + VByte 组合将二进制索引由 79.66 MB 压缩至 25.21 MB，节省 68.35\% 的空间；暖缓存条件下，布尔检索平均耗时增加约 18\%。在倒排表长度差异较大的一条查询中，平方根步长跳表使比较次数减少 89.1\%，全量指针存储约占 Raw 倒排记录的 2.8\%。基于 12 条查询、333 条人工相关性标注，原文语义编码的平均 P@10 和 MRR@10 分别为 0.533 和 1.000，均高于本实验中的 TF-IDF 与词项向量聚合方法。上述质量比较受查询规模、标注范围和编码输入差异的限制。
    \par\vspace{5pt}
  \end{minipage}
  \vspace*{8mm}
\end{titlepage}
\pagenumbering{arabic}

\section{实验设置}
\subsection{数据与任务}
课程实验包括分词与规范化（任务 1）、倒排表构建（任务 2）、位置索引与压缩（任务 3），以及布尔检索、向量空间模型和预训练模型检索（任务 4）\cite{req}。实验使用教学子集 Qilin-IR-50K 的 \path{documents.jsonl}，每篇文档包含原始编号 \path{doc_id}、标题 \texttt{title} 和正文 \texttt{content}。检索质量评价采用 12 条查询，其中 11 条为自拟查询，1 条取自 \path{queries.jsonl}；数据集查询日志未用于定量评价。

程序使用 Python、jieba、NumPy、SciPy、Sentence Transformers 和 PyTorch，模型均在 CPU 上运行。各阶段中间结果保存为 JSONL、二进制索引、CSR 稀疏矩阵或 NumPy 数组。

\subsection{实验流程与统计口径}

先在 100 篇和 1,000 篇文档上进行正确性检查，再对 50,000 篇文档运行完整实验。本文表格和结论均采用全量实验数据。其余统计约定如下：
\begin{enumerate}
  \item 文件大小采用十进制单位，即 \(1\,\mathrm{MB}=10^6\) 字节。
  \item 检索计时在词典预先加载、缓存已预热的条件下进行，以多次运行的中位数作为单条查询耗时；布尔检索每组重复 10 次。
  \item 不同模型的余弦相似度数值不直接比较，检索质量统一以相关性评价指标衡量。
  \item 代码文件与输出数据的对应关系见附录~\ref{app:files}。
\end{enumerate}


\section{文档分词与规范化（任务 1）}
\subsection{分词方法}
标题和正文分别分词，并分别记录位置；任务 3 的短语匹配不跨字段。统计词项数量时，合并两个字段的分词结果。实验采用以下两种方法：
\begin{itemize}
  \item \textbf{前向最大匹配（FMM）}：使用 jieba 内置词典的词频表，最大词长设为 6。中文片段按最长词匹配，英文、数字和符号等非中文片段按字符类别切分。
  \item \textbf{隐马尔可夫模型（HMM）}：调用 jieba 的 \texttt{finalseg} 模块，使用转移概率和发射概率进行 Viterbi 解码，不启用 jieba 的混合分词模式。
\end{itemize}

\subsection{规范化处理}
两种分词结果采用相同的规范化步骤：进行 NFKC 规范化，将英文转为小写，将纯数字词项（含小数）替换为 \texttt{<NUM>}，并去除标点及其他特殊符号。随后根据自定义的 31 词停用词表，分别生成保留停用词（Keep）与去除停用词（Drop）的词项序列，用于任务 4 的对照实验。

\subsection{结果对比与分析}
\begin{table}[H]
\centering\small
\caption{50,000 篇文档的分词统计}
\begin{tabular}{@{}lrr@{}}\toprule
指标 & FMM & HMM \\ \midrule
分词耗时 / s & 15.411 & 39.471 \\
Keep 词项出现次数 & 6,759,956 & 6,631,225 \\
Drop 词项出现次数 & 5,999,688 & 5,750,165 \\
Keep 词汇表大小 & 115,215 & 388,632 \\
Drop 词汇表大小 & 115,184 & 388,601 \\
停用词移除比例 & 11.25\% & 13.29\% \\
\bottomrule\end{tabular}
\tablesource{output/1/segmentation_stats_all.json}
\end{table}

本次测试中，HMM 的分词耗时约为 FMM 的 \(2.56\) 倍，Keep 词汇表规模约为 FMM 的 \(3.37\) 倍。以下样例展示了两种方法的切分差异：
\begin{itemize}
  \item 标题“C哥预测还准不准”中，FMM 将相关片段切为“准 / 不准”，HMM 则保留为“准不准”。
  \item 标题“手持挂烫机奥克斯”中，HMM 将“机奥克斯”合并为一个词项。FMM 分开了“机”和品牌名，但也将“挂烫”切为“挂 / 烫”。
\end{itemize}
这些样例说明，两种方法均存在边界处理不当的情况。HMM 较大的词汇表可能与不同的切分边界有关，但仅凭词汇表规模和少量样例，不能判断整体分词质量。后续倒排索引、压缩和布尔检索采用 FMM-Keep；TF-IDF 与语义检索另设 Keep / Drop 对照。

\section{倒排表构建（任务 2）}
\subsection{文档级倒排索引}
按文档输入顺序分配从 0 开始的内部编号 \(d\)，并保留其与原始 \path{doc_id} 的映射。每个词项 \(t\) 对应一个升序文档编号列表，同一文档只记录一次；列表长度即文档频率 \(df(t)\)。构建时先使用 \texttt{array(\textquotesingle I\textquotesingle)} 保存倒排表，再按词项排序输出 JSONL。

\subsection{单层跳表指针}
为减少交集运算中的逐项比较，在倒排表上等距设置单层跳跃指针。对于长度为 \(df(t)\) 的倒排表，默认步长为
\begin{equation}\label{eq:step}
  s(t)=\max\bigl(2,\lfloor\sqrt{df(t)}\rfloor\bigr),
\end{equation}
在位置 \(0,s,2s,\dots\) 处设置指向下一跳点的指针，\(df<4\) 的短表不设置跳跃边。该规则在全量索引中生成 263,858 条跳跃边；不同步长的存储和检索开销见第~\ref{sec:skip}~节。

\begin{table}[H]
\centering\small
\caption{文档级倒排索引统计}
\begin{tabular}{@{}lr@{}}\toprule
指标 & 结果 \\ \midrule
索引文档数 / 词汇量 & 50,000 / 115,215 \\
词项--文档关系数（\(\sum_t df(t)\)） & 4,074,930 \\
单层跳跃边总数（\(\sqrt{df}\) 规则） & 263,858 \\
最高文档频率词项 & “话题”（\(df=46{,}024\)） \\
文档级 JSONL 大小 / 构建耗时 & 33.62 MB / 3.43 s \\
\bottomrule\end{tabular}
\tablesource{output/2/index_stats_fmm_keep_all.json}
\end{table}

使用 \path{02_probe_index.py} 检查索引结构。全部 115,215 个词项的倒排表均严格升序且无重复，跳跃边未越界，实际步长与设定一致。本实验未实现课程选做的多层跳表，以下分析仅涉及单层等距跳表。

\section{倒排表扩展与压缩（任务 3）}
\subsection{词项位置索引}
将倒排记录扩展为三元组 \((d,P_t^{\mathrm{title}},P_t^{\mathrm{body}})\)。标题和正文分别从 0 开始记录规范化、保留停用词后的词项下标，以保留停用词所在的位置；短语匹配限定在同一字段内。位置索引采用分桶和外部排序构建，使用 64 个临时分桶，每条中间记录占 10 字节。全量数据包含 6,759,956 个词项位置，构建耗时为 19.15 s。分桶处理可减少一次性载入的记录数。

\begin{table}[H]
\centering\small
\caption{位置索引统计}
\begin{tabular}{@{}lr@{}}\toprule
指标 & 结果 \\ \midrule
词项位置出现次数 & 6,759,956（与 FMM-Keep 词项数完全一致） \\
词项--文档关系数 & 4,074,930（与任务 2 完全一致） \\
位置索引 JSONL 大小 / 构建耗时 & 82.30 MB / 19.15 s \\
\bottomrule\end{tabular}
\tablesource{output/3/positional_stats_fmm_keep_all.json}
\end{table}

使用 \path{03_check_positions.py} 逐词核对后，位置索引中的文档编号列表与任务 2 的倒排表一致，各字段的位置列表也均严格升序。

\subsection{索引压缩：Front Coding 与 Gap + VByte}
按照课程课件介绍的方法\cite{slides}，分别压缩词典和倒排记录。四组对照使用相同的位置索引，仅改变编码方式：
\begin{itemize}
  \item \textbf{Front Coding}：将排序后的词典按每块 8 个词项分组，块首保存完整词项，其余词项保存与前一词项的公共前缀字节数及后缀。
  \item \textbf{Gap + VByte}：文档编号和位置均保存相邻差值，首项相对 \(-1\) 取差，使差值为正；随后采用含 7 位数据位和 1 位终止标志的变长字节编码。Raw 对照组使用 u32 定长编码。
\end{itemize}

\begin{table}[H]
\centering\small
\caption{四种索引编码的空间占用}
\begin{tabular}{@{}lrrrrr@{}}\toprule
方案 & 词典 / MB & 倒排 / MB & 元数据 / MB & 总计 / MB & 节省率 \\ \midrule
Raw & 0.956 & 75.939 & 2.765 & 79.660 & --- \\
Front & 0.614 & 75.939 & 2.765 & 79.318 & 0.43\% \\
VByte & 0.956 & 21.830 & 2.765 & 25.551 & 67.93\% \\
Combined & 0.614 & 21.830 & 2.765 & 25.209 & 68.35\% \\
\bottomrule\end{tabular}
\tablesource{output/3/compression_summary_fmm_keep_all.json}
\end{table}

\begin{figure}[H]
\centering
\begin{tikzpicture}
\begin{axis}[
 width=0.90\linewidth,height=5.0cm,
 ybar stacked,bar width=19pt,
 symbolic x coords={Raw,Front,VByte,Combined},xtick=data,
 ymin=0,ymax=87,ytick={0,20,40,60,80},
 ylabel={空间占用 / MB},
 ylabel style={font=\small},
 tick label style={font=\small},
 ymajorgrids,grid style={black!15},
 axis line style={black!50},axis x line*=bottom,axis y line*=left,
 legend cell align=left,
 legend style={column sep=6pt,font=\footnotesize,draw=none,at={(0.5,1.02)},anchor=south,legend columns=3},
 enlarge x limits=0.17,
 ]
 \addplot+[color=black,draw=black!75,fill=black!70] coordinates {(Raw,75.939)(Front,75.939)(VByte,21.830)(Combined,21.830)};
 \addplot+[color=black,draw=black!75,fill=black!25] coordinates {(Raw,2.765)(Front,2.765)(VByte,2.765)(Combined,2.765)};
 \addplot+[color=black,draw=black!75,fill=white] coordinates {(Raw,0.956)(Front,0.614)(VByte,0.956)(Combined,0.614)};
 \legend{倒排记录,元数据,词典}
\end{axis}
\end{tikzpicture}
\caption{四种方案的二进制索引空间构成}
\end{figure}

Front Coding 使词典体积减少 \(35.78\%\)。由于词典仅占原索引的约 1.2\%，总空间减少 \(0.43\%\)。Gap + VByte 将倒排记录从 75.94 MB 压缩至 21.83 MB，减少 \(71.25\%\)，是组合方案节省空间的主要来源。

上述比例均以字段相同的二进制 Raw 索引（79,659,905 B）为基准，不包含 JSONL 转为二进制格式带来的体积变化。总大小包括词典、倒排记录和索引元数据，不包括共享文档映射。通过 \path{03_verify_compression.py} 校验，Combined 模式解压后的 115,215 个词项、4,074,930 条倒排记录和 6,759,956 个位置均与原索引一致；其余模式各抽样 100 个词项，结果一致。

\section{信息检索方法比较（任务 4）}
\subsection{布尔检索（任务 4A）}
\subsubsection{复杂布尔表达式与处理顺序}
设计 9 条布尔查询，覆盖 AND、OR、NOT、嵌套括号和纯 NOT 表达式，所用词项均在 FMM 词汇表中。基准测试采用 Raw 索引，按表达式的原始顺序处理，并关闭跳表。

\begin{table}[H]
\centering\small
\caption{9 条布尔查询的检索结果}
\begin{tabularx}{\linewidth}{@{}Xrrr@{}}\toprule
查询表达式 & 命中数 & 中位耗时 / ms & 比较次数 \\ \midrule
\textsf{((登机牌 AND 习惯) OR 收藏) AND NOT 话题} & 46 & 134.77 & 5,409 \\
\textsf{(平价 OR 护肤品) AND (学生 AND 党)} & 157 & 4.76 & 4,052 \\
\textsf{(山西 AND 壁画) OR (古建 AND 艺术)} & 2 & 1.34 & 787 \\
\textsf{(旅行 AND 攻略) AND NOT 广告} & 1,016 & 19.22 & 54,253 \\
\textsf{(考研 AND (英语 OR 政治)) AND 经验} & 20 & 5.56 & 4,076 \\
\textsf{(北京 AND 故宫) OR (西安 AND 兵马俑)} & 26 & 1.64 & 996 \\
\textsf{(咖啡 AND (做法 OR 教程)) AND NOT 广告} & 19 & 14.37 & 51,579 \\
\textsf{NOT (话题 OR 分享)} & 3,749 & 177.63 & 46,240 \\
\textsf{(分享 AND 攻略) AND 旅行}（高频词优先） & 182 & 22.16 & 14,084 \\
\bottomrule\end{tabularx}
\tablesource{output/4/boolean_benchmark_fmm_all.json}
\end{table}

\noindent\textbf{处理顺序。}分别按原始左结合顺序和基于 \(df\) 的估计顺序执行查询；后者优先处理 AND 链中预计结果较小的子表达式。本次实验中，重排前后的比较次数基本一致。例如，高频词优先的查询分别需要 14,084 和 14,096 次比较，差异小于 \(0.1\%\)，未观察到稳定的耗时改善。

线性双指针归并的单次交集最坏时间复杂度为 \(O(|A|+|B|)\)。执行顺序仍会影响中间结果的长度，但本次查询集未体现出明显收益，不能据此认定重排对所有查询均无效。含高频词“话题”（\(df=46{,}024\)）和纯 NOT 补集运算的查询耗时较高，倒排解码及补集处理可能是主要开销。

\subsubsection{压缩前后检索效率}
四种编码方式使用相同的 9 条查询，均保持原始处理顺序并关闭跳表。下表先以每条查询的重复运行中位数作为单条耗时，再计算 9 条查询间的平均值和中位数。

\begin{table}[H]
\centering\small
\caption{四种索引编码的布尔检索效率（9 条查询汇总）}
\begin{tabular}{@{}lrrrr@{}}\toprule
方案 & 平均耗时 / ms & 中位耗时 / ms & 词典加载 / s & 索引大小 / MB \\ \midrule
Raw & 42.38 & 14.37 & 0.164 & 79.66 \\
Front & 42.70 & 14.23 & 0.195 & 79.32 \\
VByte & 49.49 & 15.04 & 0.186 & 25.55 \\
Combined & 50.07 & 15.00 & 0.196 & 25.21 \\
\bottomrule\end{tabular}
\end{table}

Combined 的空间节省率为 68.35\%，暖缓存条件下的平均耗时由 42.38 ms 增至 50.07 ms，增加约 \(18\%\)。这表明在本实验实现中，压缩并未缩短检索时间，可能与逐字节解码的 CPU 开销有关。Front 与 Raw 的查询耗时接近。冷缓存、部分倒排解码等条件下的表现尚未测试，不能由当前结果推断。

\subsubsection{短语检索与位置索引的效果}
对 5 个短语进行位置检索，要求词项在同一字段内按顺序相邻，并与相同词项的普通 AND 查询比较。过滤率定义为 \(1-\text{短语命中数}/\text{AND 命中数}\)。

\begin{table}[H]
\centering\small
\caption{短语检索与普通 AND 查询的结果对照}
\begin{tabular}{@{}lrrrrr@{}}\toprule
短语 & 短语命中 & AND 命中 & 过滤率 & Raw / ms & Combined / ms \\ \midrule
收藏 登机牌 & 1 & 1 & --- & 1.24 & 1.41 \\
平价 护肤品 & 2 & 13 & 84.6\% & 1.67 & 1.84 \\
北京 故宫 & 2 & 23 & 91.3\% & 1.41 & 1.74 \\
减肥 运动 & 4 & 187 & 97.9\% & 3.71 & 4.41 \\
旅行 攻略 & 264 & 1,018 & 74.1\% & 9.62 & 11.40 \\
\bottomrule\end{tabular}
\end{table}

以“减肥 运动”为例，普通 AND 查询返回 187 篇包含两个词项的文档，短语检索保留其中 4 篇，过滤率为 \(97.9\%\)。两者的差异来自是否要求词项在同一字段内相邻。“收藏 登机牌”在两种方式下均命中 1 篇，说明这条查询未产生额外过滤。Combined 的短语检索耗时均高于 Raw；按表中数据计算，增幅约为 10--23\%。

\subsubsection{跳表步长对存储与检索的影响}\label{sec:skip}
选择以下 3 条 AND 查询，比较无跳表及 5 种跳表步长。每条跳跃边保存目标文档号和倒排内偏移，均为 u32，因此按每条 8 B 估算指针存储开销。
\begin{itemize}
  \item Q1：\texttt{(}话题 AND 登机牌\texttt{) AND} 习惯；
  \item Q2：\texttt{(}攻略 AND 旅行\texttt{) AND} 民宿；
  \item Q3：\texttt{(}考研 AND 英语\texttt{) AND} 复习。
\end{itemize}

\begin{table}[H]
\centering\small
\caption{不同跳表步长下的比较次数与中位耗时（Raw 索引）}
\begin{tabular}{@{}lrrrrrr@{}}\toprule
查询 & 无跳表 & 2 & 4 & 8 & 16 & $\sqrt{df}$ \\ \midrule
Q1 比较次数 & 46,189 & 23,194 & 11,717 & 6,083 & 3,445 & 5,018 \\
中位耗时 / ms & 160.2 & 139.2 & 134.9 & 131.8 & 131.0 & 136.5 \\
Q2 比较次数 & 5,707 & 6,838 & 6,324 & 6,145 & 5,983 & 5,810 \\
中位耗时 / ms & 8.97 & 9.26 & 9.33 & 9.39 & 9.30 & 9.27 \\
Q3 比较次数 & 1,576 & 1,875 & 1,785 & 1,738 & 1,664 & 1,648 \\
中位耗时 / ms & 2.77 & 2.83 & 2.88 & 2.88 & 2.75 & 2.85 \\
\bottomrule\end{tabular}
\end{table}

\begin{table}[H]
\centering\small
\caption{不同跳表步长的指针存储开销}
\begin{tabular}{@{}lrrr@{}}\toprule
步长 & 跳跃边总数 & 指针存储 / MB & 占 Raw 倒排记录比例 \\ \midrule
2 & 1,949,530 & 15.60 & 20.5\% \\
4 & 962,007 & 7.70 & 10.1\% \\
8 & 467,852 & 3.74 & 4.9\% \\
16 & 224,561 & 1.80 & 2.4\% \\
$\sqrt{df}$ & 263,858 & 2.11 & 2.8\% \\
\bottomrule\end{tabular}
\tablesource{output/4/skip_storage_fmm_all.json}
\end{table}

\begin{figure}[H]\centering
\begin{tikzpicture}
\begin{axis}[
 width=0.90\linewidth,height=5.5cm,
 ymode=log,
 line width=0.8pt,mark size=2pt,
 symbolic x coords={none,s2,s4,s8,s16,sqrt},xtick=data,
 xticklabels={无跳表,2,4,8,16,$\sqrt{df}$},
 ymin=1000,ymax=100000,
 ylabel={比较次数（对数轴）},ylabel style={font=\small},
 tick label style={font=\small},
 ymajorgrids,grid style={black!15},axis line style={black!50},axis x line*=bottom,axis y line*=left,
 enlarge x limits=0.10,
 legend cell align=left,
 legend style={column sep=6pt,font=\footnotesize,draw=none,at={(0.5,1.04)},anchor=south,legend columns=3},
 ]
 \addplot+[color=black,draw=black,mark=*,mark options={fill=black},solid] coordinates {(none,46189)(s2,23194)(s4,11717)(s8,6083)(s16,3445)(sqrt,5018)};
 \addplot+[color=black,draw=black!65,mark=square*,mark options={fill=black!65,solid},dashed] coordinates {(none,5707)(s2,6838)(s4,6324)(s8,6145)(s16,5983)(sqrt,5810)};
 \addplot+[color=black,draw=black!75,mark=triangle*,mark options={fill=white,solid},densely dotted] coordinates {(none,1576)(s2,1875)(s4,1785)(s8,1738)(s16,1664)(sqrt,1648)};
 \legend{Q1,Q2,Q3}
\end{axis}
\end{tikzpicture}
\caption{三条查询的比较次数随跳表步长变化（对数轴）}
\end{figure}

Q1 中，“话题”的文档频率为 46,024，“登机牌”为 21，倒排表长度差异较大。平方根步长使比较次数减少 \(89.1\%\)，中位耗时由 160.2 ms 降至 136.5 ms，减少约 \(14.8\%\)；固定步长 16 的比较次数减少 \(92.5\%\)，耗时降至 131.0 ms。时间降幅小于比较次数降幅，与跳表未减少倒排解码量有关。

Q2 和 Q3 中，跳表未带来稳定的耗时改善。固定步长 2 反而使比较次数分别增加约 \(19.8\%\) 和 \(19.0\%\)，因为判断能否跳跃也需要比较。平方根步长的比较次数增幅分别约为 \(1.8\%\) 和 \(4.6\%\)。

平方根步长的指针估算大小为 2.11 MB，占 Raw 倒排记录的 2.8\%，低于固定步长 2 的 15.60 MB。在这 3 条查询上，它能减少 Q1 的比较次数，对 Q2、Q3 的影响较小；固定步长 16 则使用更少的指针，并在 Q1 上取得更低耗时。因此，当前结果不足以确定普遍最优的步长。

\subsection{向量空间模型（任务 4B）}
采用原始词频与平滑 IDF 计算权重，文档和查询使用相同规则，标题与正文同权。文档矩阵采用 \texttt{float32} 类型的 CSR 稀疏存储，相关性由余弦相似度计算：
\begin{equation}
 w_{t,d}=tf(t,d)\left[\ln\frac{N+1}{df(t)+1}+1\right], \qquad
 \operatorname{sim}(q,d)=\frac{\mathbf q^\mathsf T\mathbf d}{\|\mathbf q\|_2\|\mathbf d\|_2}.
\end{equation}

\begin{table}[H]\centering\small
\caption{TF-IDF 模型的构建统计}
\begin{tabular}{@{}lrrr@{}}\toprule
指标 & Keep & Drop & 下降比例 \\ \midrule
词汇表规模 & 115,215 & 115,184 & 0.027\% \\
词项出现次数 & 6,759,956 & 5,999,688 & 11.25\% \\
矩阵非零元素 & 4,074,930 & 3,800,308 & 6.74\% \\
CSR 文件 / MB & 22.73 & 21.40 & 5.85\% \\
构建耗时 / s & 9.41 & 8.56 & 9.06\% \\
\bottomrule\end{tabular}
\end{table}

删除 31 个停用词后，词项出现次数减少约 76 万。使用与语义检索相同的 12 条查询比较 Keep 和 Drop 模型；候选数指与查询具有正相似度的文档数，Jaccard@10 为两个 Top-10 文档集合的交集大小与并集大小之比。

\begin{table}[H]\centering\small
\caption{TF-IDF 保留与去除停用词的检索结果}
\begin{tabular}{@{}lrrrc@{}}\toprule
查询 & Keep 候选 & Drop 候选 & Top-10 重合 & Jaccard@10 \\ \midrule
收藏 登机牌 & 821 & 821 & 10/10 & 1.000 \\
有 收藏 登机牌 的 习惯 吗 & 42,849 & 18,817 & 10/10 & 1.000 \\
山西 古建 壁画 & 86 & 86 & 10/10 & 1.000 \\
平价 护肤品 学生党 & 2,518 & 2,518 & 9/10 & 0.818 \\
旅行 攻略 民宿 & 4,555 & 4,555 & 10/10 & 1.000 \\
考研 英语 复习 经验 & 2,887 & 2,887 & 10/10 & 1.000 \\
咖啡 做法 教程 & 2,476 & 2,476 & 9/10 & 0.818 \\
北京 故宫 门票 & 1,463 & 1,463 & 10/10 & 1.000 \\
减肥 运动 食谱 & 2,179 & 2,179 & 10/10 & 1.000 \\
衣柜 颜色 搭配 & 3,196 & 3,196 & 10/10 & 1.000 \\
酒酿好处（真实查询） & 339 & 339 & 10/10 & 1.000 \\
适合学生党的平价护肤品…（长查询） & 43,048 & 23,351 & 10/10 & 1.000 \\
\bottomrule\end{tabular}
\tablesource{output/4/tfidf_comparison/}
\end{table}

12 条查询中，有 10 条的 Keep / Drop Top-10 集合相同，另 2 条各有 9 篇重合。对含功能词的第 2、12 条查询，去除停用词后，候选数分别从 42,849 降至 18,817、从 43,048 降至 23,351，但 Top-10 集合未变。本次结果表明，停用词处理对这两条查询的候选规模影响较大，对高排名集合的影响较小；集合相同并不表示内部排序完全一致。

\subsection{语义检索（任务 4C）}
\subsubsection{方法与规模}
使用 \path{BAAI/bge-small-zh-v1.5} 生成 512 维向量，最大输入长度为 512 个模型 token，在 CPU 上推理。Raw 直接编码“标题 + 正文”原文；Seg-Keep 对 FMM 词项逐词独立编码，再按词频加权平均并进行 L2 归一化；Seg-Drop 在聚合前去除停用词。三种方法均使用模型自带的 tokenizer，另设 TF-IDF 作为词汇匹配基线。

\begin{table}[H]\centering\small
\caption{50,000 篇文档的语义向量规模与记录耗时}
\begin{tabular}{@{}lrrrr@{}}\toprule
方法 & 文档数 & 维度 & 向量文件 / MB & 记录耗时 / s \\ \midrule
Raw & 50,000 & 512 & 102.40 & 1,277.9 \\
Seg-Keep & 50,000 & 512 & 102.40 & 8.5\textsuperscript{*} \\
Seg-Drop & 50,000 & 512 & 102.40 & 10.5\textsuperscript{*} \\
\bottomrule\end{tabular}
\vspace{2pt}\parbox{\linewidth}{\footnotesize\textsuperscript{*} Seg 计时对应词项向量缓存（115,215 词独立编码）建立后的文档聚合与写盘，未计入首次词项编码成本；不可与 Raw 的完整推理时间直接比较。}
\end{table}

Raw 与 Seg 的差异不只涉及预先分词：Raw 的输入超过 512 个模型 token 时会被截断，而 Seg 聚合使用全部规范化词项；Seg 还替换了数字并删除标点。此外，逐词编码后平均会改变文本的上下文表示。因此，这一对照不能单独用于评价预先分词的影响。

\subsubsection{分词与停用词对结果集合的影响}
\begin{table}[H]\centering\small
\caption{语义检索方法的 Top-10 集合重合程度（Jaccard@10）}
\begin{tabularx}{\linewidth}{@{}Xcc@{}}\toprule
查询 & Raw / Seg-Keep & Seg-Keep / Seg-Drop \\ \midrule
收藏 登机牌 & 0.111 & 0.429 \\
有 收藏 登机牌 的 习惯 吗 & 0.000 & 0.000 \\
山西 古建 壁画 & 0.053 & 0.538 \\
平价 护肤品 学生党 & 0.053 & 0.667 \\
旅行 攻略 民宿 & 0.053 & 0.667 \\
考研 英语 复习 经验 & 0.053 & 0.538 \\
咖啡 做法 教程 & 0.333 & 0.538 \\
北京 故宫 门票 & 0.111 & 1.000 \\
减肥 运动 食谱 & 0.000 & 0.818 \\
衣柜 颜色 搭配 & 0.111 & 0.538 \\
酒酿好处 & 0.053 & 0.333 \\
适合学生党的平价护肤品…（长查询） & 0.053 & 0.333 \\
\midrule
\textbf{平均} & 0.082 & 0.492 \\
\bottomrule\end{tabularx}
\tablesource{output/4/semantic/comparison_all_*.json}
\end{table}

Raw 与 Seg-Keep 的平均 Jaccard@10 为 0.082，说明两种方法返回的 Top-10 集合差异较大。逐词独立编码后求平均不保留词序，也不直接表示上下文中的词项搭配，这可能是差异的原因之一；输入截断和规范化差异也可能产生影响。Seg-Keep 与 Seg-Drop 的平均 Jaccard@10 为 0.492，低于 TF-IDF Keep / Drop 的集合重合程度，说明本次测试中停用词对词项向量聚合结果的影响更大。

\subsubsection{人工相关性评价（P@10 / MRR@10）}
将 12 条查询、4 种方法的 Top-10 结果按“查询--文档”对去重，共得到 333 条标注，其中相关 98 条、不相关 235 条。标注模板由脚本生成，再根据标题和正文预览逐条判断相关性，标签取 0 或 1。评价指标为：
\begin{equation}
 P@10(q)=\frac{\#\{\text{前 10 篇中标为相关的文档}\}}{10},\qquad
 \operatorname{MRR}@10=\frac{1}{|Q|}\sum_{q\in Q}\frac{1}{r_q},
\end{equation}
其中 \(r_q\) 为查询 \(q\) 首个相关文档的排名，前 10 篇无相关结果时令 \(1/r_q=0\)。

\begin{table}[H]\centering\small
\caption{各查询的 P@10 与汇总指标}
\begin{tabularx}{\linewidth}{@{}Xcccc@{}}\toprule
查询 & Raw & Seg-Keep & Seg-Drop & TF-IDF \\ \midrule
收藏 登机牌 & 0.10 & 0.10 & 0.10 & 0.10 \\
有 收藏 登机牌 的 习惯 吗 & 0.10 & 0.00 & 0.10 & 0.10 \\
山西 古建 壁画 & 0.20 & 0.10 & 0.20 & 0.20 \\
平价 护肤品 学生党 & 0.70 & 0.10 & 0.10 & 0.00 \\
旅行 攻略 民宿 & 1.00 & 0.90 & 0.80 & 0.50 \\
考研 英语 复习 经验 & 0.80 & 0.40 & 0.30 & 0.40 \\
咖啡 做法 教程 & 0.80 & 0.80 & 0.70 & 0.60 \\
北京 故宫 门票 & 0.30 & 0.00 & 0.00 & 0.10 \\
减肥 运动 食谱 & 1.00 & 0.90 & 0.90 & 0.80 \\
衣柜 颜色 搭配 & 0.60 & 0.20 & 0.20 & 0.10 \\
酒酿好处 & 0.40 & 0.10 & 0.10 & 0.40 \\
适合学生党的平价护肤品…（长查询） & 0.40 & 0.00 & 0.10 & 0.00 \\
\midrule
\textbf{平均 P@10} & 0.533 & 0.300 & 0.300 & 0.275 \\
\textbf{MRR@10} & \textbf{1.000} & 0.561 & 0.771 & 0.750 \\
\bottomrule\end{tabularx}
\tablesource{output/4/semantic/precision_all_*.json}
\end{table}

\begin{figure}[H]\centering
\begin{tikzpicture}
\begin{axis}[
 width=0.90\linewidth,height=5.0cm,
 ybar,bar width=12pt,
 symbolic x coords={Raw,SegKeep,SegDrop,TFIDF},xtick=data,
 xticklabels={Raw,Seg-Keep,Seg-Drop,TF-IDF},
 ymin=0,ymax=1.15,ytick={0,0.2,0.4,0.6,0.8,1.0},
 ylabel={评价指标},ylabel style={font=\small},
 tick label style={font=\small},
 ymajorgrids,grid style={black!15},axis line style={black!50},axis x line*=bottom,axis y line*=left,
 enlarge x limits=0.17,
 legend cell align=left,
 legend style={column sep=6pt,font=\footnotesize,draw=none,at={(0.5,1.02)},anchor=south,legend columns=2},
 nodes near coords,nodes near coords style={font=\scriptsize,/pgf/number format/fixed,/pgf/number format/precision=3},
 ]
 \addplot+[color=black,draw=black!75,fill=white] coordinates {(Raw,0.533)(SegKeep,0.300)(SegDrop,0.300)(TFIDF,0.275)};
 \addplot+[color=black,draw=black!75,fill=black!25] coordinates {(Raw,1.000)(SegKeep,0.561)(SegDrop,0.771)(TFIDF,0.750)};
 \legend{平均 P@10,MRR@10}
\end{axis}
\end{tikzpicture}
\caption{四种检索方法的平均 P@10 与 MRR@10（12 条查询）}
\end{figure}


Raw 在本次查询集上的平均 P@10 为 0.533，MRR@10 为 1.000，即每条查询的首位结果均被标为相关。这两项汇总指标高于另外三种方法。Seg-Keep 与 Seg-Drop 的平均 P@10 均为 0.300，MRR@10 则分别为 0.561 和 0.771，表明去除停用词后首个相关结果的平均排名有所改善，但前 10 篇的平均相关比例未提高。TF-IDF 在“减肥 运动 食谱”和“咖啡 做法 教程”上的 P@10 分别为 0.80 和 0.60，在“平价 护肤品 学生党”上为 0。


\subsubsection{案例分析}
\begin{itemize}
  \item \textbf{“平价 护肤品 学生党”}：Raw 的 Top-10 中有 7 篇相关，Seg-Keep 和 Seg-Drop 各有 1 篇，TF-IDF 为 0 篇。该查询同时涉及价格、品类和适用人群，仅匹配部分词项的文档未必满足完整需求。
  \item \textbf{“北京 故宫 门票”}：Raw 的 P@10 为 0.30，检出的《逛故宫，去之前的准备工作》包含预约与门票信息；Seg-Keep 的 P@10 为 0。这一结果与逐词聚合难以保留上下文关联的特点相符，但不能据此排除其他输入差异的影响。
  \item \textbf{“减肥 运动 食谱”与“咖啡 做法 教程”}：四种方法的 P@10 均在 0.60--1.00 之间。查询词能较直接地描述文档主题，TF-IDF 在这两个案例中也能获得较多相关结果。
  \item \textbf{“酒酿好处”}：该查询取自数据集。Raw 与 TF-IDF 的 P@10 均为 0.40，Seg-Keep 与 Seg-Drop 均为 0.10。单个案例尚不足以说明造成差异的具体原因。
\end{itemize}

\section{讨论与结论}
\subsection{主要结论}
\begin{enumerate}
  \item \textbf{分词。}本次运行中，HMM 耗时约为 FMM 的 2.56 倍，Keep 词汇表规模约为 FMM 的 3.37 倍。两种方法均存在切分不当的样例，尚未进行系统的分词准确率评价。
  \item \textbf{索引与压缩。}文档级索引和位置索引包含相同的 4,074,930 条词项--文档关系。Combined 模式经全量解压校验，空间节省率为 68.35\%，其中倒排记录的 Gap + VByte 编码贡献了主要压缩量。
  \item \textbf{布尔检索。}AND 链重排在本次查询集上未带来明显收益。倒排表长度差异较大的 Q1 中，平方根步长跳表使比较次数减少 89.1\%，耗时减少约 14.8\%；另外两条查询未表现出稳定的耗时改善。Combined 在暖缓存条件下的平均检索耗时较 Raw 增加约 18\%。
  \item \textbf{短语检索。}位置条件可排除不满足字段内连续匹配的文档。5 条查询中，1 条与普通 AND 的结果相同，其余 4 条的过滤率为 74.1--97.9\%。
  \item \textbf{检索质量。}基于 12 条查询和 333 条标注，Raw 的平均 P@10 与 MRR@10 分别为 0.533 和 1.000，高于另外三种方法。停用词处理对 TF-IDF 的 Top-10 集合影响较小，对词项聚合结果及 MRR@10 的影响较大。
\end{enumerate}

\subsection{限制}
本实验有以下限制：
\begin{enumerate}
  \item 查询仅有 12 条，且部分主题相近，结果不能代表更广泛的检索需求。
  \item 标注池仅覆盖四种方法的 Top-10 并集，无法估计全集召回率。相关性由一人根据标题和正文预览判断，未进行标注者间一致性检验。
  \item Raw 与 Seg 在输入截断、规范化和聚合方式上存在差异，不能将质量变化完全归因于分词；Seg 计时还未包含首次词项编码成本。
  \item 检索计时采用暖缓存条件，尚未测试冷缓存和按需解码，也未对更大规模查询集验证跳表步长的效果。
\end{enumerate}
后续可扩大真实查询样本，并将 \path{search_logs} 中的点击记录作为补充相关性信号，同时保留人工判断，以进一步检查结论的适用范围。

\vspace{2mm}
\begin{thebibliography}{9}
\setlength{\itemsep}{0pt}
\bibitem{req}《实验一：信息检索部分》，课程实验要求，2026。
\bibitem{slides}《Web 信息处理与应用》，第三节“网页索引”课程课件，索引压缩部分。
\bibitem{qilin} THUIR / 小红书，Qilin 数据集及课程教学子集 Qilin-IR-50K；\url{https://huggingface.co/datasets/THUIR/Qilin}。
\bibitem{logs} 本实验各阶段输出的 JSON 统计、布尔检索基准、语义检索对照 JSON 与人工相关性 CSV（随代码和索引归档）。
\end{thebibliography}

\Needspace{10\baselineskip}
\appendix
\section{代码与输出数据对照}\label{app:files}
全部脚本在项目根目录执行，命令形式为 \texttt{python src/}\textnormal{脚本名}\texttt{.py}。100 篇和 1,000 篇规模用于正确性检查；\texttt{all} 表示 50,000 篇全量数据。各任务的代码与主要输出如下，具体命令见 \path{src/README.md}。

\begin{table}[H]\centering\small
\caption{实验任务、代码文件与输出数据对照}
\begin{tabularx}{\linewidth}{@{}lY Y@{}}\toprule
任务 & 代码文件 & 主要输出（相对 output/） \\ \midrule
数据体检 & \path{00_check_data.py} & \path{data_check.json} \\
任务 1 & \path{01_segment.py}、\path{01_compare_samples.py} & \path{1/tokens_*.jsonl}、\path{1/segmentation_*} \\
任务 2 & \path{02_build_index.py}、\path{02_probe_index.py} & \path{2/inverted_*.jsonl}、\path{2/index_stats_*.json} \\
任务 3 & \path{03_build_positions.py}、\path{03_compress_index.py}、\path{stage3b_codec.py}、\path{03_check_positions.py}、\path{03_verify_compression.py} & \path{3/positional_*.jsonl}、\path{3/compressed_*/}、\path{3/compression_*} \\
任务 4A & \path{04A_boolean_retrieval.py}、\path{04A_skip_storage.py} & \path{4/boolean_benchmark_*.json}、\path{4/skip_storage_*.json} \\
任务 4B & \path{04B_build_tfidf.py}、\path{04B_search_tfidf*.py} & \path{4/tfidf_fmm_*.npz} 等、\path{4/tfidf_comparison/} \\
任务 4C & \path{04C_build/compare/evaluate_semantic.py} & \path{4/semantic/}（向量、对照、标注 CSV、P@10） \\
\bottomrule\end{tabularx}
\end{table}

\vspace{2mm}

\noindent\textbf{提交前核对：}补全组长姓名、学号和科大云盘链接，核对已有成员信息；复核 \path{output/4/semantic/relevance_labels_*.csv} 中的相关性标签，尤其是新增 8 条查询对应的 212 条标题与正文预览判断记录；将源代码、核心索引和统计文件按课程要求归档。

\end{document}
